\phantomsection
\section*{Lời mở đầu}
\addcontentsline{toc}{section}{Lời mở đầu}
\vspace{10pt}
\fontsize{14pt}{21pt}\selectfont

Trong phạm vi bài toán mô phỏng, cán bộ cần tra cứu ai đang cư trú ở địa chỉ nào, thuộc hộ nào, theo diện thường trú hay tạm trú, và ghi nhận các lần chuyển đến, chuyển đi hoặc chuyển sang hộ khác. Nhóm giả định dữ liệu ban đầu được theo dõi bằng sổ sách hoặc các tệp Excel rời rạc; cách tổ chức này có thể gây trùng lặp, khó tra cứu lịch sử và khó xác định người đã sửa dữ liệu.

Từ bài toán trên, Nhóm 38 lựa chọn đề tài \textbf{Xây dựng cơ sở dữ liệu quản lý cư dân và biến động cư trú tại một phường mô phỏng}. Mô hình lưu riêng các đợt cư trú theo thời gian và các biến động có ghi tài khoản, thời điểm nhập, giúp tra cứu lịch sử đã ghi nhận. Việc cập nhật được kiểm soát bằng quyền và trigger; thiết kế hiện tại chưa bảo đảm lịch sử bất biến và chưa có nhật ký kiểm toán đầy đủ. Nhóm cài đặt cơ sở dữ liệu trên hệ quản trị PostgreSQL 17, dùng thêm ràng buộc loại trừ, trigger và hàm PL/pgSQL cho những quy tắc mà khóa và CHECK thông thường không diễn đạt được.

Báo cáo gồm sáu chương, bám theo yêu cầu của bài tập lớn. Chương 1 trình bày bài toán thực tiễn, các quy tắc nghiệp vụ và sơ đồ phân rã chức năng của hệ thống. Chương 2 xây dựng mô hình thực thể liên kết. Chương 3 ánh xạ mô hình đó sang mô hình dữ liệu quan hệ, mô tả cấu trúc các bảng và các ràng buộc dữ liệu. Chương 4 và Chương 5 lần lượt trình bày mười câu truy vấn bằng đại số quan hệ và mười câu truy vấn bằng SQL. Chương 6 tìm hiểu vấn đề bảo mật trong hệ cơ sở dữ liệu, đưa ra ba ví dụ minh họa trên chính cơ sở dữ liệu của nhóm và đánh giá các giới hạn còn lại. Phần Kết luận tổng kết kết quả, giới hạn và hướng phát triển. Trang cuối là bảng phân công công việc và đánh giá mức độ đóng góp của các thành viên.

Toàn bộ dữ liệu dùng để chạy thử trong báo cáo là bộ dữ liệu giả lập 636 bản ghi do nhóm tự sinh, không phải thông tin của người dân thật.
